%------------------------------------------------------------------------------%

\usepackage{integracao_e_series}
\DeclareMathOperator{\bessel}{J}

\title{Séries de Potências -- Manipulações}

%------------------------------------------------------------------------------%
\begin{document}

\addtitle

%------------------------------------------------------------------------------%
\section{Função Definida por Série de Potências}

%------------------------------------------------------------------------------%
\begin{frame}{Função Definida por Série de Potências}
  Se
  \[
    \sum_{n=0}^\infty c_n (x-a)^n
  \]
  \pause
  possui um raio de convergência \(R>0\),
  \pause
  então a função
  \[
    f(x)=\sum_{n=0}^\infty c_n(x-a)^n
  \]
  \pause
  está definida no intervalo $(a-R,\,a+R)$ \\[3mm] \pause
  é infinitamente derivável nesse intervalo
\end{frame}

%------------------------------------------------------------------------------%
\begin{frame}{Manipulando Funções Definidas por Séries}
  Podemos fazer mudanças de variáveis e somar as séries

  \pause
  \vspace{2\baselineskip}
  Desde que estejamos atentos ao intervalo de convergência
\end{frame}

%------------------------------------------------------------------------------%
\addtocounter{exemplo}{1}
\begin{frame}{Exemplo \theexemplo}
  Escreva \(f(x) = \frac{1}{1-x^2}\) como uma série de potências

  \pause
  \vspace{\baselineskip}
  Sabemos que
  \[
    \frac{1}{1-u}
    \;=\; \sum_{n=0}^\infty u^n
    \;=\; 1 + u + u^2 + u^3 + \cdots + u^n + \cdots
  \]
  para \(\abs{u}<1\)
\end{frame}

%------------------------------------------------------------------------------%
\begin{frame}{Exemplo \theexemplo}
  \onslide<+->{Fazendo \(u = x^2\)}\onslide<+->{, temos
  \[
    \frac{1}{1-u}
    \;=\; \sum_{n=0}^\infty u^n
    \;=\; 1 + u + u^2 + u^3 + \cdots + u^n + \cdots
  \]}
  \[
    \onslide<+->{\frac{1}{1-x^2}
    \;=\; \sum_{n=0}^\infty (x^2)^n}
    \onslide<+->{\;=\; \sum_{n=0}^\infty x^{2n}}
    \onslide<+->{\;=\; 1 + x^2 + x^4 + x^6 + \cdots + x^{2n} + \cdots}
  \]

  \onslide<+->{Analisando o intervalo de convergência}
  \begin{align*}
    \onslide<+->{\abs{u} & < 1        \\}
    \onslide<+->{x^2     & < 1        \\}
    \onslide<+->{\abs{x} & < \sqrt{1} \\}
    \onslide<+->{-1 < x  & < 1}
  \end{align*}
\end{frame}

%------------------------------------------------------------------------------%
\addtocounter{exemplo}{1}
\begin{frame}{Exemplo \theexemplo}
  Escreva \(f(x) = \frac{1}{x+2}\) como uma série de potências

  \pause
  \vspace{\baselineskip}
  Sabemos que
  \[
    \frac{1}{1-u}
    \;=\; \sum_{n=0}^\infty u^n
    \;=\; 1 + u + u^2 + u^3 + \cdots + u^n + \cdots
  \]
  para \(\abs{u}<1\)
\end{frame}

%------------------------------------------------------------------------------%
\begin{frame}{Exemplo \theexemplo}
  Precisamos manipular $f$
  \pause
  \[
    f(x)
    \;=\; \frac{1}{x+2}
    \pause
    \;=\; \frac{1}{2\left(1+\dfrac{x}{2}\right)}
    \pause
    \;=\; \frac{1}{2}\,\frac{1}{\left[1-\left(-\dfrac{x}{2}\right)\right]}
  \]

  \pause
  Fazendo \(u=-\frac{x}{2}\)\pause, temos
  \[
    f(x)
    \;=\; \frac{1}{2} \, \sum_{n=0}^\infty \left(-\frac{x}{2}\right)^n
    \pause
    \;=\; \sum_{n=0}^\infty \frac{1}{2} \, \frac{(-1)^n\,x^n}{2^n}
    \pause
    \;=\; \sum_{n=0}^\infty \frac{(-1)^n}{2^{n+1}}\,x^n
  \]
\end{frame}

%------------------------------------------------------------------------------%
\begin{frame}{Exemplo \theexemplo}
  \onslide<+->{Analisando o intervalo de convergência}
  \begin{align*}
    \onslide<+->{\abs{u}    \;\;        & < 1 \\[1mm]}
    \onslide<+->{\abs{-\frac{x}{2}}     & < 1 \\[1mm]}
    \onslide<+->{\frac{\abs{x}}{2} \;\; & < 1 \\[1mm]}
    \onslide<+->{\abs{x} \;\;           & < 2}
  \end{align*}
  \onslide<+->{A série converge para \(x\in(-2, 2)\)}
\end{frame}

%------------------------------------------------------------------------------%
\section{Derivada da Série}

%------------------------------------------------------------------------------%
\begin{frame}{Derivação Termo a Termo}
  Podemos calcular a derivada de $f(x)$ construindo a série
  \begin{align*}
    \action<+->{f'(x) & = {\color{blue}\frac{d}{dx} \Bigg(} \sum_{n=0}^\infty c_n(x-a)^n {\color{blue}\Bigg)} \\[2mm]}
    \action<+->{      & = \sum_{n=0}^\infty {\color{blue}\frac{d}{dx} \Big(} c_n(x-a)^n {\color{blue}\Big)} \\[2mm]}
    \action<+->{      & = \sum_{{\color<5->{red}n=1}}^\infty c_n \,n\, \big(x-a\big)^{n-1}}
    \action<5->{\qquad\text{Note que o índice dessa série começa em {\color{red}1}}}
  \end{align*}
  \action<+->{que converge em $(a-R, a+R)$}
\end{frame}

%------------------------------------------------------------------------------%
\addtocounter{exemplo}{1}
\begin{frame}{Exemplo \theexemplo}
  Calcule a derivada da função de Bessel
  \pause
  \[
  \bessel_0(x) = \sum_{n=0}^\infty \frac{\,(-1)^n\,x^{2n}\,}{2^{2n}\,(n!)^2}
  \]
\end{frame}

%------------------------------------------------------------------------------%
\begin{frame}{Exemplo \theexemplo}
  \begin{columns}
    \column{0.5\baselineskip}
    \begin{align*}
      \onslide<+->{\bessel'_0(x)}
      \onslide<+->{& = \frac{d}{dx} \left(\sum_{n=0}^\infty \frac{(-1)^nx^{2n}}{2^{2n}(n!)^2}\right) \\[3mm]}
      \onslide<+->{& = \sum_{n=0}^\infty \frac{d}{dx} \left(\frac{(-1)^nx^{2n}}{2^{2n}(n!)^2}\right) \\[3mm]}
      \onslide<+->{& = \sum_{n=0}^\infty \frac{(-1)^n}{2^{2n}(n!)^2} \frac{d}{dx} \left(x^{2n}\right)}
    \end{align*}
    \column{0.5\baselineskip}
    \begin{align*}
      \onslide<+->{\bessel'_0(x)
        & = \sum_{n=1}^\infty \frac{(-1)^n}{2^{2n}(n!)^2} \; 2n x^{2n-1} \\[3mm]}
      \onslide<+->{& = \sum_{n=1}^\infty \frac{\,(-1)^n\, n\, x^{2n-1}\,}{2^{2n-1}(n!)^2}}
    \end{align*}
  \end{columns}
\end{frame}

%------------------------------------------------------------------------------%
\section{Integral da Série}

%------------------------------------------------------------------------------%
\begin{frame}{Integração Termo a Termo}
  Para calcularmos a integral de $f$ construímos a série
  \begin{align*}
    \action<+->{\int f(x)dx & = {\color{blue}\int} \left(\sum_{n=0}^\infty c_n(x-a)^n\right){\color{blue}dx} \\[2mm]}
    \action<+->{            & = \sum_{n=0}^\infty \left({\color{blue}\int} c_n(x-a)^n {\color{blue}dx}\right) \\[2mm]}
    \action<+->{            & = \sum_{n=0}^\infty c_n\frac{(x-a)^{n+1}}{n+1}+ C}
\end{align*}
  \action<+->{que também converge em $(a-R, a+R)$}
\end{frame}

%------------------------------------------------------------------------------%
\addtocounter{exemplo}{1}
\begin{frame}{Exemplo \theexemplo}
  Encontre uma série de potências centrada em \(x=0\) para a função
  \[
    f(x) = \ln \left( 1 + x^2 \right)
  \]
\end{frame}

%------------------------------------------------------------------------------%
\begin{frame}{Exemplo \theexemplo}
  Ainda não temos uma técnica para obter a série diretamente\pause,
  mas sabemos que
  \[
    \frac{d}{dx} \ln\left( 1+x^2 \right)
    \;=\; \frac{2x}{1+x^2}
    \pause\;=\; \frac{2x}{1-\left(-x^2\right)}
  \]

  \pause
  \vspace{\baselineskip}
  Enquanto que a soma da Série Geométrica com \(\alpha=2x\) e \( r = -x^2 \) é
  \[
    \sum_{n=0}^\infty 2x \left(-x^2\right)^n
    \;=\; \frac{2x}{1-\left(-x^2\right)}
  \]
\end{frame}

%------------------------------------------------------------------------------%
\begin{frame}{Exemplo \theexemplo}
  Nos pontos onde a série converge
  \[
    \frac{d}{dx}\ln\left( 1+x^2 \right)
    \pause\;=\; \frac{2x}{1-\left(-x^2\right)}
    \pause\;=\; \sum_{n=0}^\infty 2x \left(-x^2 \right)^n
    \pause\;=\; 2 \sum_{n=0}^\infty (-1)^n\,x^{2n+1}
  \]

  \pause
  Basta integrar para obtermos a série desejada
\end{frame}

%------------------------------------------------------------------------------%
\begin{frame}{Exemplo \theexemplo}
  \begin{align*}
    \onslide<+->{\ln\left( 1+x^2 \right)
      = \int \frac{d}{dx} \ln\left( 1+x^2 \right) dx}
    \onslide<+->{& = \int 2 \sum_{n=0}^\infty (-1)^n\,x^{2n+1} dx \\[2mm]}
    \onslide<+->{& = 2 \sum_{n=0}^\infty (-1)^n \int x^{2n+1} dx \\[2mm]}
    \onslide<+->{& = 2 \sum_{n=0}^\infty (-1)^n \left(\frac{x^{2n+2}}{\,2n+2\,} + c_n\right) \\[2mm]}
    \onslide<+->{& = C + \sum_{n=0}^\infty (-1)^n \frac{x^{2n+2}}{\,n+1\,}}
  \end{align*}
\end{frame}

%------------------------------------------------------------------------------%
\begin{frame}{Exemplo \theexemplo}
  Falta determinar \(C\)

  \pause
  Avaliando a função e a série em $x=0$
  \pause
  \[
    \ln\left( 1+x^2 \right) \;=\; \ln(1) \;=\; 0
  \]
  \pause
  \vspace{-\baselineskip}
  \[
    C + \sum_{n=0}^\infty (-1)^n \frac{x^{2n+2}}{\,n+1\,}
    = C
  \]
  \pause
  portanto \(C=0\)\pause, então
  \[
    \ln\left( 1+x^2 \right) = \sum_{n=0}^\infty \frac{(-1)^n}{n+1}\, x^{2n+2}
  \]
\end{frame}

%------------------------------------------------------------------------------%
\begin{frame}{\contentsname}
  \tableofcontents
\end{frame}

%------------------------------------------------------------------------------%
\end{document}
%------------------------------------------------------------------------------%
